\begin{afterword}
\summaryhead

The post compares two heroes who storm a room held by an armed kidnapper. One is an invulnerable superhero saving two hundred children; the other is a police officer saving three adults. Who is the greater hero, and who gets the comic book?

By the halo effect, the author argues, strength and invulnerability make a character seem braver too. Fame, the author suggests, adds to every other trait: Gandhi, partly protected by celebrity, is remembered as the model of nonviolence, while anonymous marchers who were beaten are forgotten.

The author then separates two senses of ``greater.'' As an outcome, saving two hundred is better than saving three, and anyone forced to choose should save the two hundred. As revealed virtue, someone who would risk death for three shows more courage than someone who would do it only for two hundred. Choosing the riskier or smaller rescue in order to look virtuous is a moral failure, and virtuous people look for safer ways to save more. Whatever risk remains shows real courage, so the unpowered officer reveals more virtue than Superman.

\responsehead

The second half of the post is careful moral reasoning. The distinction between the better outcome and the greater revealed virtue is stated precisely, the comparison between someone who would risk death for three and someone who would do it only for two hundred holds as stated, and the warning against choosing the smaller rescue to display courage follows from it. The conclusion about the officer follows from the post's own definitions.

The first half names a bias it does not show. The claim that invulnerable heroes ``seem to possess more of the heroic traits of courage'' is inferred from the halo effect, with no study of how people judge such characters; the support offered is a line from a comic. The halo effect itself is stated as ``perceptions of all positive traits are correlated,'' a stronger form than the evidence discussed under ``The Halo Effect'' allows.

The claim about fame is offered with ``seems,'' and the author is unsure whether it is original; no evidence is given. Its example is one-sided. Gandhi's celebrity did spare him some of what his followers suffered: at Dharasana marchers were clubbed while he was in jail. But his prominence also made him a target, of a mob in 1897 and of several attempts on his life before the one that killed him. The post counts the protection and not the risk. Its later softening, ``somewhat-but-not-entirely protected,'' still ranks his courage below that of anonymous marchers without comparing their risks.

In short: a sound distinction between good outcomes and revealed courage, together with a bias that is named but not shown, and with a claim about Gandhi's risk that counts only one side.
\end{afterword}
